% Problem-section file following ../_template.tex.
\section{Statistical strength of CGLMP measurements}
\label{sec:problem-25}

\paragraph{Problem.}
For every $d\geq3$, do the standard CGLMP Fourier--phase measurements maximize
the relative-entropy statistical strength against local realism among all
projective $d$-outcome measurements on the fixed state
$\lvert\Phi_d\rangle=d^{-1/2}\sum_{j=0}^{d-1}\lvert j,j\rangle$, when the
setting distribution is also optimized?  For a behavior $p(a,b\mid x,y)$, a
distribution $\mu(x,y)$ on the four setting pairs, and the local polytope
$\mathcal L$, define
\begin{equation}
  S(p;\mu)
  :=\inf_{\ell\in\mathcal L}
    \sum_{a,b,x,y}\mu(x,y)p(a,b\mid x,y)
    \log_2\!\frac{p(a,b\mid x,y)}{\ell(a,b\mid x,y)},
  \label{eq:p25-statistical-strength}
\end{equation}
Set $S^\star(p):=\sup_{\mu\in\Delta(\{0,1\}^2)}S(p;\mu)$; by
Eq.~\eqref{eq:p25-statistical-strength}, this is the optimized asymptotic
evidence rate against the best local model.  The candidate bases are
\begin{equation}
  \begin{aligned}
    \lvert a;x\rangle
      &=\frac{1}{\sqrt d}\sum_{j=0}^{d-1}
        \exp\!\left(\frac{2\pi i}{d}j(a+\alpha_x)\right)\lvert j\rangle,\\
    \lvert b;y\rangle
      &=\frac{1}{\sqrt d}\sum_{j=0}^{d-1}
        \exp\!\left(\frac{2\pi i}{d}j(-b+\beta_y)\right)\lvert j\rangle,
  \end{aligned}
  \label{eq:p25-cglmp-bases}
\end{equation}
where $\alpha_0=0$, $\alpha_1=-1/2$, $\beta_0=1/4$, and $\beta_1=3/4$.
Let $p_{\mathrm{CGLMP}}$ denote the behavior produced on
$\lvert\Phi_d\rangle$ by the bases in Eq.~\eqref{eq:p25-cglmp-bases}, and
write $p_M$ for the behavior produced by any other measurement choice $M$ on
$\lvert\Phi_d\rangle$.  The conjectured optimality of
Eq.~\eqref{eq:p25-cglmp-bases} is
\begin{equation}
  S^\star(p_{\mathrm{CGLMP}})
  =\sup_M S^\star(p_M),
  \label{eq:p25-cglmp-statistical-optimum}
\end{equation}
where the supremum is over two projective $d$-outcome measurements per party.
Equation~\eqref{eq:p25-cglmp-statistical-optimum} is the question to be
resolved.

\subsection*{Status}

Unsolved

\subsection*{Source}

Gill explicitly proposes the global statistical-strength optimality of the
CGLMP measurement construction for a fixed maximally entangled state
\sourcecite{ref:p25-gill}{Gil07}.

\subsection*{Progress}

\begin{itemize}
  \item Van Dam, Gill, and Grünwald established the operational interpretation
  of Eq.~\eqref{eq:p25-statistical-strength} as the asymptotic evidence rate of
  a Bell experiment and formulated the joint optimization problem
  \sourcecite{ref:p25-statistical-strength}{DGG05}.

  \item Acín, Gill, and Gisin numerically found the CGLMP measurement bases for
  their relative-entropy searches at $d=3$ and $d=4$.  Their globally best
  state was not maximally entangled, so this does not resolve the fixed-state
  equality in Eq.~\eqref{eq:p25-cglmp-statistical-optimum}
  \sourcecite{ref:p25-acin}{AGG05}.

  \item Gill reports that numerical searches with the maximally entangled
  state fixed found only the CGLMP measurements for both Euclidean and
  relative-entropy criteria, but treats optimality as conjectural; the search
  does not prove Eq.~\eqref{eq:p25-cglmp-statistical-optimum} for arbitrary
  $d$
  \sourcecite{ref:p25-gill}{Gil07}.
\end{itemize}

\subsection*{References}

\begin{enumerate}[label={},leftmargin=4.8em,labelsep=0.5em]
  \footnotesize
  \item[\textup{[DGG05]}]\label{ref:p25-statistical-strength}
  W. van Dam, R. D. Gill, and P. D. Grünwald,
  ``The Statistical Strength of Nonlocality Proofs,''
  \emph{IEEE Transactions on Information Theory} \textbf{51}, 2812--2835
  (2005).
  \href{https://doi.org/10.1109/TIT.2005.851738}{doi:10.1109/TIT.2005.851738};
  \href{https://arxiv.org/abs/quant-ph/0307125}{arXiv:quant-ph/0307125}.

  \item[\textup{[AGG05]}]\label{ref:p25-acin}
  A. Acín, R. Gill, and N. Gisin,
  ``Optimal Bell Tests Do Not Require Maximally Entangled States,''
  \emph{Physical Review Letters} \textbf{95}, 210402 (2005).
  \href{https://doi.org/10.1103/PhysRevLett.95.210402}{doi:10.1103/PhysRevLett.95.210402};
  \href{https://arxiv.org/abs/quant-ph/0506225}{arXiv:quant-ph/0506225}.

  \item[\textup{[Gil07]}]\label{ref:p25-gill}
  R. D. Gill, ``Better Bell Inequalities (Passion at a Distance),'' in
  \emph{Asymptotics: Particles, Processes and Inverse Problems}, IMS Lecture
  Notes--Monograph Series \textbf{55}, 135--148 (2007).
  \href{https://doi.org/10.1214/074921707000000328}{doi:10.1214/074921707000000328};
  \href{https://arxiv.org/abs/math/0610115}{arXiv:math/0610115}.
\end{enumerate}

\subsection*{Comment}

The remaining gap is a global proof or counterexample to
Eq.~\eqref{eq:p25-cglmp-statistical-optimum} for arbitrary $d$.  The shared
state is fixed; this differs from joint state--measurement optimization and
from maximizing the linear CGLMP violation in Problem~24.

\subsection*{Tag}

Bell nonlocality; Convex optimization; Measurement theory; Qudit systems;
Quantum foundations

\subsection*{ID}

\texttt{op\_a34f0e2d6489068f}
